% ===========================================================================
% Diagramme de cas d'utilisation — acteurs et cas principaux.
%
% Deux modèles d'acteurs coexistent et le rapport doit l'assumer (conflit
% C-02) : le cahier des charges v1.0 en nomme quatre (visiteur, candidat,
% responsable RH, administrateur) ; l'implémentation en porte quatre AUTRES —
% ADMIN, RH, RECRUTEUR, CANDIDAT.
%
% Sens de l'héritage vérifié dans backend/app/core/rbac_policy.csv :
%   g, RH, RECRUTEUR   → le RH hérite des permissions du recruteur
% La politique porte aussi `g, ADMIN, RH`, volontairement NON représentée :
% une généralisation UML modélise un rôle métier, pas un niveau de droits.
% L'administrateur est un acteur TECHNIQUE — en faire un RH spécialisé lui
% prêterait un rôle dans les décisions d'embauche et l'accès aux salaires.
% Côté Casbin son périmètre reste total, mais par le joker `p, ADMIN, *, *`,
% qui rend d'ailleurs la ligne `g, ADMIN, RH` sans effet.
%
% Les cas sont ORDONNÉS PAR ACTEUR : un ordre thématique obligeait chaque
% acteur interne à traverser toute la figure pour joindre ses cas.
% ===========================================================================
\begin{tikzpicture}[node distance=4mm]

  % ═══ Cas d'utilisation — colonne centrale, groupés par acteur ═══════════
  \node[pqbox, text width=40mm] (uc1)
    {Consulter et rechercher\\les offres publiées};
  \node[pqbox, below=of uc1, text width=40mm] (uc2)
    {Créer un compte\\et s'authentifier};
  \node[pqbox, below=of uc2, text width=40mm] (uc3)
    {Déposer un CV\\et postuler};
  \node[pqbox, below=of uc3, text width=40mm] (uc4)
    {Suivre ses candidatures\\et leur analyse};
  \node[pqbox, below=of uc4, text width=40mm] (uc5)
    {Confirmer un créneau\\d'entretien};
  \node[pqbox, below=of uc5, text width=40mm] (uc6)
    {Passer le test technique};
  \node[pqbox, below=of uc6, text width=40mm] (uc7)
    {Négocier et signer\\le contrat};
  \node[pqbox, below=of uc7, text width=40mm] (uc8)
    {Consulter le classement\\des candidats};
  \node[pqbox, below=of uc8, text width=40mm] (uc9)
    {Planifier un entretien\\et assigner un recruteur};
  \node[pqbox, below=of uc9, text width=40mm] (uc10)
    {Évaluer un entretien\\et décider};
  \node[pqbox, below=of uc10, text width=40mm] (uc11)
    {Publier et gérer\\une offre};
  \node[pqbox, below=of uc11, text width=40mm] (uc12)
    {Éditer un contrat\\et activer l'employé};
  \node[pqbox, below=of uc12, text width=40mm] (uc13)
    {Piloter budgets,\\employés et carrières};
  \node[pqbox, below=of uc13, text width=40mm] (uc14)
    {Administrer comptes,\\rôles et journal d'audit};

  % ═══ Frontière du système ═══════════════════════════════════════════════
  \begin{scope}[on background layer]
    \node[pqboundary, fit=(uc1)(uc14), inner sep=8pt] (sysbox) {};
  \end{scope}
  \node[pqcaption, anchor=south] at (sysbox.north) {PIQBIT};

  % ═══ Acteurs externes — placés au niveau de leurs cas ═══════════════════
  \node[pqactor, left=16mm of uc1, text width=24mm]  (visit) {Visiteur};
  \node[pqactor, left=16mm of uc5, text width=24mm]  (cand)  {Candidat};
  \node[pqactor, right=16mm of uc9, text width=24mm] (recr)  {Recruteur};
  \node[pqactor, right=16mm of uc12, text width=24mm] (rh)   {Responsable RH};
  \node[pqactor, right=16mm of uc14, text width=24mm] (adm)  {Administrateur};

  % ═══ Héritage de permissions (fait réel de la politique Casbin) ═════════
  % La flèche va du rôle HÉRITIER vers le rôle hérité : « le RH hérite du
  % recruteur ». L'inverse se lirait comme une délégation, qui n'existe pas.
  \draw[pqflow, draw=piqbitExt] (rh.north) --
    node[pqcard, right, text=piqbitExt, align=left] {hérite de} (recr.south);

  % ═══ Associations acteur → cas ══════════════════════════════════════════
  \draw[pqrel] (visit.east) -- (uc1.west);
  \draw[pqrel] (visit.east) -- (uc2.west);

  \draw[pqrel] (cand.east) -- (uc3.west);
  \draw[pqrel] (cand.east) -- (uc4.west);
  \draw[pqrel] (cand.east) -- (uc5.west);
  \draw[pqrel] (cand.east) -- (uc6.west);
  \draw[pqrel] (cand.east) -- (uc7.west);

  \draw[pqrel] (recr.west) -- (uc8.east);
  \draw[pqrel] (recr.west) -- (uc9.east);
  \draw[pqrel] (recr.west) -- (uc10.east);

  \draw[pqrel] (rh.west) -- (uc11.east);
  \draw[pqrel] (rh.west) -- (uc12.east);
  \draw[pqrel] (rh.west) -- (uc13.east);

  \draw[pqrel] (adm.west) -- (uc14.east);

  % ═══ Note sur le conflit d'acteurs ══════════════════════════════════════
  \node[pqwarn, below=7mm of sysbox.south, anchor=north, text width=120mm]
    {\textbf{Deux modèles d'acteurs à réconcilier (conflit C-02).}
     Le cahier des charges v1.0 nomme visiteur, candidat, responsable RH et
     administrateur. L'implémentation porte \textsc{admin}, \textsc{rh},
     \textsc{recruteur} et \textsc{candidat}. Le « visiteur » n'est pas un
     rôle stocké mais l'\emph{absence} de compte ; le rôle \textsc{recruteur}
     est apparu après la rédaction du cahier des charges. Le RH peut tout ce
     que peut le recruteur, et la généralisation est justifiée : les deux
     jouent le même rôle dans le cycle de recrutement, le RH y ajoutant les
     salaires et l'interne. L'administrateur en revanche est un acteur
     \emph{technique} : la politique lui accorde un périmètre total
     (\texttt{p, ADMIN, *, *}), mais le modéliser comme un RH spécialisé lui
     prêterait un rôle dans les décisions d'embauche et l'accès aux
     rémunérations, ce qui romprait la séparation des devoirs.};

\end{tikzpicture}
